\section{Task Definition}
\label{sec:task}
% Paragraph role P01: 输入输出
Given a current speaker audio segment $a_s$ and the listener's first frame $I_l$, the model generates listener video $V_l$ and audio $a_l$ for the corresponding interval. The image provides the listener's appearance, and the audio provides information for producing feedback. The intended use case involves one or two short utterances.
\cnnote{统计待补}{训练与测试中实际片段长度及句数分布需单独报告；不将“一两句”的应用设定写成已经统计确认的数据分布。}

% Paragraph role P02: 倾听／发声边界
We focus on brief feedback while the other participant holds the conversational floor. A response may consist only of visual movement, or may include an acknowledgment or laughter. Quiet listening is also within scope: neither sound nor salient motion is required in every clip. Listening role and vocal activity are distinct; a short vocalization alone does not establish that an utterance is a backchannel.

% Paragraph role P03: 训练推理信息边界
At inference, no dialogue history, future listener frames, target audio, or ground-truth reaction labels are provided. Paired listener media and annotations can be used for training. In particular, guided training uses annotated reaction descriptions, whereas inference uses descriptions predicted by Qwen, as detailed in \cref{sec:method}.

% Paragraph role P04: 一对多属性及限制
The current utterance does not uniquely determine all listener behavior. Different facial responses, sounds, and timings can be appropriate. We treat each observed response as a training example without assuming it is the only valid output. This task property does not itself establish that the model learns controllable diversity.
